\scanpage{172}
\mspage{30}

\inserted{\uline{Cor.} 2.14} Notons quand même qu'\deleted{un espace topologique} \inserted{si $\unsure{\mathscr X|\mathcal X}$ est de la} forme $\operatorname{Top}(T)$, $T$ un espace topologique, alors $\mathscr X$ est loc. tot. $0$-connexe ssi il est loc irréductible i.e. ssi $T$ est loc. irréductible. On est ramené à prouver que cette dernière condition, implique \deleted{\unsure{alors}} la condition e) de 2.6, \deleted{\unsure{et}} on peut prendre pour $S$ l'ensemble \inserted{ordonné des} ouverts connexes non vides de $S$, elle devient alors évidente.

Ceci dit, il faut se rappeler que la théorie cohomologique des espaces topologiques irréductibles n'est nullement triviale -- comme nous le montre FAC de Serre ! D'autre part, c'est Serre qui \unsure{m'avait} fait \unsure{observer} que pour des coefficients \uline{constants}, la cohomologie d'un espace irréductible est triviale.

Dans le cas actuel, prouvons déjà, \deleted{\unsure{que}} \inserted{si $\mathscr X$ est $0$-connexe et loc. tot. $0$-connexe}, \inserted{que} \unsure{tt|tout} \unsure{revêtement} de \deleted{\unsure{de} $\mathscr X$ est} l'objet final de $\mathscr X$ est trivial. Il suffit de voir que si $U$ est $0$-connexe, \deleted{il est} $U\longrightarrow e$ est un iso. \deleted{On} Mais considérons $U\hookrightarrow U\times U$, on sait que c'est un \unsure{sommande|sous-objet} \unsure{direct|discret}, et comme $U\times U$ est connexe, c'est un iso. Donc $U\longrightarrow e$ mono, \deleted{et comme} donc \unsure{sommande|sous-objet} \unsure{direct|discret}, et comme $U\ne\varnothing$, on a $U\longrightarrow e$ iso,

Reste \unsure{à|,} prouver que pour \unsure{tt|le} faisceau

% archive sheet number: 172
